\section{Discussion and Limitations}
\label{sec:discussion}

\medskip\noindent\textbf{Interpreting the trace-guided routing signals.}
The two gates in Failure-Trace Allocator (FTA) instantiate a selective-deferral
framework in which the frontier answer is retained by default and overridden only
on runtime-identifiable subsets without any access to gold labels.
The Low-Depth Guard (LDG) addresses a structural weakness: when the frontier search
terminates with a \textit{single weak branch}---signalled by the \texttt{override\_reason}
field---the resulting answer is drawn from a narrow slice of the candidate space.
In such cases, the external majority vote is empirically more reliable.
The External-Consensus Fallback (ECO) targets a different pattern: when the
direct-agreement trigger fires---signalling that the frontier selected its own answer
confidently---yet all three independent external baselines unanimously prefer a different
answer, that unanimous disagreement suggests that the frontier answer may be less reliable in this evaluation.
Both signals are available as logged metadata from normal frontier execution;
no additional inference or labelling is required.

\medskip\noindent\textbf{Relation to pooled aggregation and selective routing.}
The pooled four-answer ensemble is an important comparator because classical ensemble results suggest that simple aggregation can already be powerful when component errors are not perfectly aligned.
In this manuscript, the pooled baseline is deliberately treated as a strong reference point rather than as a weak baseline to be dismissed.
Failure-Trace Allocator is competitive with that pooled baseline and remains ahead by point estimate on both Final-300 and Aggregate-720, but the paired/bootstrap confidence intervals include zero.
Accordingly, the claim is not that trace-guided routing is broadly superior to ensembling.
The narrower claim is that interpretable, gold-free answer-source routing can
help in this matched-budget answer-selection protocol and remain competitive with
a simple pooled ensemble built from the same already-generated answers.

\medskip\noindent\textbf{Additional variants and why they were not adopted.}
A TALE-default router produced zero switches on seeds~41 and~61,
equalling TALE performance exactly on both splits (80.33\% and 54.17\%,
respectively).
On seed~71 (Final-300), that router made one net positive switch over TALE
(78.67\% vs.\ 78.33\%), a difference of a single example that is negligible
compared to the $+8.33$~pp advantage of Failure-Trace Allocator (FTA) over TALE on the same split.
The TALE-default design assumed a reliable agreement region that did not materialise
consistently across evaluation seeds, making the router effectively inactive in most
conditions.

An extra-action variant proposed collecting extra frontier or TALE-retry actions for predicted
failure cases.
The first 40-case overlapping pilot showed nonnegative action-value signal, but
an independent 80-case follow-up (disjoint from the pilot) yielded
net $-1$ for extra frontier actions and net $-4$ for extra TALE retries relative to Failure-Trace Allocator (FTA),
with meaningful regressions on the corresponding held-out subsets (8.9\% for frontier, 11.1\% for TALE retry).
The failure to replicate on disjoint data is the key indicator: the pilot signal
appears to have been artefact-specific rather than generalisable.

A cluster-level selector showed only limited positive movement in offline replay,
with a best-rule local gain of $+0.67$~pp on final-300, but it did not provide
aggregate evidence or a robust policy improvement and was therefore not adopted for the final policy.
A robust parser variant produced zero accuracy-changing switches on
final-300 and aggregate-720.
Taken together, these outcomes indicate that the targeted residual patterns are either
rare or too weakly actionable in the current evaluation setting.

\medskip\noindent\textbf{How aggregate-720 changes the evidence posture.}
The most substantive evidence shift in this work is the transition from a single-run
300-example evaluation (seed 41, where the CI lower bound vs.\ L1 still crossed zero
at $[-0.48, +5.71]$) to the disjoint Aggregate-720 evidence (seeds 41, 61, 71).
By combining three independently sampled unbiased evaluation splits, the stratified
bootstrap CI narrows sufficiently that the lower bound vs.\ every external baseline
is strictly positive.
The source-stratified design controls for
seed-specific variance, and the zero-overlap check on example-id and question-hash
confirms independence.
The result is that the point-estimate advantage of Failure-Trace Allocator (FTA)
is consistent across all three seeds, which is the primary reason the aggregate
CIs no longer include zero. This strengthens the case for the selective-deferral
framing in the tested setting, but it does not transform the paper into a broad
claim of universal routing superiority.

\medskip\noindent\textbf{Evidence scope and generalization.}
The main findings are scoped to the tested evaluation setting.
What this evidence does \emph{not} establish is generalization beyond that scope:
whether the same failure-trace signals are informative on other providers,
benchmarks, or budget levels is an open empirical question.
Existing repository artifacts include earlier MATH-500, AIME, OpenAI/Cohere, and
budget-sweep outputs, but these do not contain the same four-method per-example
records and frontier trace metadata needed to replay Failure-Trace Allocator (FTA) under the current
protocol.
They are therefore unsuitable as a clean second benchmark/provider/budget result
for this selective-deferral study without new matched-budget generation.
The most efficient next replication would be a second GSM8K provider run or a
same-provider second benchmark with the identical frontier, L1, S1, and TALE
logging contract.
Section~\ref{sec:limitations} enumerates the scope boundaries explicitly.
